\documentclass[11pt]{article}
\usepackage{amsmath}
\usepackage[hidelinks]{hyperref}
\title{What the native MRWS fault runs can establish}
\date{}
\begin{document}
\maketitle

The MultiRobot Warehouse Scheduling simulator can witness a simulated
robot leaving a cell and a station receiving inventory. It cannot
establish custody of a unique physical item. Its fault campaign also
includes runs in which a robot tries to move into an occupied cell
and the simulator stops before applying the move. Omitting those
runs would turn a failed safety-and-service comparison into a biased
recovery test.

\paragraph{Source and run boundary.}
We ran \href{https://github.com/automaticdai/MRWS_MultiRobot_Warehouse_Scheduling/tree/941338e9ea4c4eaebb87f8374f80fd576497ad01}
{the pinned MRWS source} on its \(7\times11\) layout, with the native
\texttt{simple-interrupt} scheduler, four robots from the layout,
capacity three, ten inventory classes, native order generation and
the fault-aware flag on. The development-only protocol chose seeds
0--15, a horizon of 240 ticks, and actuator/sensor fault probabilities
\(1/64\) and \(1/128\) per applicable draw. Battery faults were off.
These are test inputs, not estimates of field failure frequencies.
The Python hash seed was fixed. Genetic mode was not run: a local
import substitute raises if that optional scheduler is invoked.
No simulator source file was edited. An in-process recorder replaced
the dashboard sender without network I/O. The move wrapper records a
native guard's inputs and then calls it unchanged. A second complete
run reproduced the captured JSONL byte for byte.

\paragraph{Results and what the logs mean.}
All 16 nominal instances completed, with 64 order completion
events. With the two enabled fault modes, one of 16 instances
completed, one reached the tick limit, and 14 stopped at the
native occupied-cell guard. All 14 rejected moves were attempted
by robots marked sensor-faulted. The guard precedes the location
update: the collision-labelled exception is a rejected simulated
move, not a recorded state with two robots in one cell. Attained
end-of-tick configurations had distinct positions. Across the
fault runs, 63 native fault events occurred with nonempty carrier
counts; the harness counts 92 later-fault/loaded-predecessor pairs
before the earlier carrier count emptied. Those are pairs, not
92 separate missions or 92 verified overlapping recovery duties.
Loaded status is reconstructed from native gain, lost and clear
events, not from task reassignment.

The native source adds a shelf's existing item class on each
pickup; it does not deplete individual stock. A station receives
the robot's inventory and the scheduler matches counts of names,
the assigned station and the assigned robot. Each recorded
order completion had a station gain in the same simulated tick.
Acknowledging a reassignment is not an order completion;
an actual source receipt is available for the counted service.
Conversely, no observation here says which unique item was delivered,
nor does the dashboard's complete view represent a robot's
knowledge during a communication break.

The runs were paired by seed only, not by a pre-generated common
trajectory of every exogenous decision. Conditional fault draws
and control actions can change later random consumption. Thus
the counts are campaign outcomes, not a matched causal estimate
of a controller improvement or fault effect size. The meaningful
same-run result is the native precondition rejection and the
distinction between fault, reschedule, move, carried inventory
and receipt. No new controller was tested.

\paragraph{Decision for a transfer.}
Use this source if the held-out obligation concerns occupied space,
class-count service, or an actual order receipt, not exclusive
custody of a unique item. Retain stops and timeouts as maximal
outcomes, and freeze allowed observations before selecting an
unseen compound case. Counting terminal guard rejection as safe,
productive recovery would be invalid. A recovery comparison needs
a real source for timely obstruction evidence and a lawful action
that handles rejection without discarding live work. Either pin a
shared exogenous workload and fault trace before running alternative
policies or limit the claims to observations on individual native
runs. These development cases do not consume a holdout or establish
I1 or I2 acceptance.
\end{document}

